\setchoices{label=(\Alph*)}

\begin{question}[index=1]
设 $\displaystyle I_R=\oint_L\frac{y\dd x-x\dd y}{(x^2+y^2+xy)^2}$，
其中 $L$ 为 $x^2+y^2=\dfrac{1}{R^2}$，方向为逆时针方向，则 $R\to0^+$ 时，
\begin{choices}[columns=2]
\item $I_R$ 是 $R$ 的 $1$ 阶无穷小．
\item $I_R$ 是 $R$ 的 $2$ 阶无穷小．
\item $I_R$ 是 $R$ 的 $3$ 阶无穷小．
\item $I_R$ 至少是 $R$ 的 $3$ 阶无穷小．
\end{choices}
\end{question}
\QuestionBreak

\begin{question}[index=2]
设函数 $f(x)$ 在 $[-\pi,\pi]$ 上连续，当
$\displaystyle F(a)=\int_{-\pi}^{\pi}[f(x)-a\cos nx]^2\dd x$（$n\in\symbf{N}_+$）取得极小值时，$a=$
\begin{choices}[columns=2,label-pos=left]
\item $\displaystyle\int_{-\pi}^{\pi} f(x)\cos nx\dd x$．
\item $\displaystyle\frac{1}{\pi}\int_{-\pi}^{\pi} f(x)\cos nx\dd x$．
\item $\displaystyle\frac{2}{\pi}\int_{-\pi}^{\pi} f(x)\cos nx\dd x$．
\item $\displaystyle\frac{1}{2\pi}\int_{-\pi}^{\pi} f(x)\cos nx\dd x$．
\end{choices}
\end{question}
\QuestionBreak

\begin{question}[index=3]
设 $p>0$，$\displaystyle\int_2^{+\infty}\frac{\mathrm dx}{(\ln x)^{(\ln x)^p}}$ 收敛，则 $p$ 的取值范围是
\begin{choices}[columns=2]
\item $0<p\le1$．
\item $p\ge2$．
\item $p\ge1$．
\item $1\le p<2$．
\end{choices}
\end{question}
\QuestionBreak

\begin{question}[index=4]
设 $L$ 是依次连接点 $(0,0)$，$(0,2)$，$(2,1)$，$(1,0)$，$(0,0)$ 的折线，
$D$ 是由 $L$ 围成的平面单连通区域，若 $P(x,y),Q(x,y)$ 均具有一阶连续偏导数，且
\[
\oint_L Q(x,y)\dd x+P(x,y)\dd y=\iint_D f(x,y)\dd\sigma,
\]
则 $f(x,y)=$
\begin{choices}[columns=4,label-pos=left]
\item $\dfrac{\partial Q}{\partial x}-\dfrac{\partial P}{\partial y}$．
\item $\dfrac{\partial Q}{\partial y}-\dfrac{\partial P}{\partial x}$．
\item $\dfrac{\partial P}{\partial y}-\dfrac{\partial Q}{\partial x}$．
\item $\dfrac{\partial P}{\partial x}-\dfrac{\partial Q}{\partial y}$．
\end{choices}
\end{question}
\QuestionBreak

\begin{question}[index=5]
设 $3$ 阶非零矩阵 $\symbfit{A}$ 满足 $A_{ij}+a_{ij}=0$（$i,j=1,2,3$），
其中 $A_{ij}$ 为元素 $a_{ij}$ 的代数余子式，则下列结论不正确的是
\begin{choices}[columns=2]
\item $(\symbfit{A}-\symbfit{E})\symbfit{x}=\symbf{0}$ 有非零解．
\item $(\symbfit{A}+\symbfit{E})\symbfit{x}=\symbf{0}$ 有非零解．
\item $(\symbfit{A}+\symbfit{E})^2\symbfit{x}=\symbf{0}$ 有非零解．
\item $(\symbfit{A}^2-\symbfit{E})\symbfit{x}=\symbf{0}$ 有非零解．
\end{choices}
\end{question}
\QuestionBreak

\begin{question}[index=6]
设 $\symbfit{A}$ 为 $n$ 阶实对称矩阵，$\symbfit{B}$ 为 $n$ 阶可逆矩阵，$\symbfit{Q}$ 为 $n$ 阶正交矩阵，
则下列矩阵中一定与 $\symbfit{A}$ 有相同的特征值的是
\begin{choices}[columns=2]
\item $\symbfit{B}^{-1}\symbfit{Q}^{\mathrm T}\symbfit{A}\symbfit{Q}\symbfit{B}$．
\item $(\symbfit{B}^{-1})^{\mathrm T}\symbfit{Q}^{\mathrm T}\symbfit{A}\symbfit{Q}\symbfit{B}^{-1}$．
\item $\symbfit{B}^{\mathrm T}\symbfit{Q}^{\mathrm T}\symbfit{A}\symbfit{Q}\symbfit{B}$．
\item $\symbfit{B}\symbfit{Q}^{\mathrm T}\symbfit{A}\symbfit{Q}(\symbfit{B}^{\mathrm T})^{-1}$．
\end{choices}
\end{question}
\QuestionBreak

\begin{question}[index=7]
若向量组 $\symbfit{\alpha}_1=(1,1,a)^{\mathrm T}$，$\symbfit{\alpha}_2=(1,a,1)^{\mathrm T}$，
$\symbfit{\alpha}_3=(a,1,1)^{\mathrm T}$ 可由向量组 $\symbfit{\beta}_1=(1,1,a)^{\mathrm T}$，
$\symbfit{\beta}_2=(-2,a,4)^{\mathrm T}$，$\symbfit{\beta}_3=(-2,a,a)^{\mathrm T}$ 线性表示，
但向量组 $\symbfit{\beta}_1,\symbfit{\beta}_2,\symbfit{\beta}_3$ 不能由向量组
$\symbfit{\alpha}_1,\symbfit{\alpha}_2,\symbfit{\alpha}_3$ 线性表示，则 $a=$
\begin{choices}[columns=4]
\item $-2$．
\item $2$．
\item $-1$．
\item $1$．
\end{choices}
\end{question}
\QuestionBreak

\begin{question}[index=8]
设随机变量 $X_1,X_2,\ldots,X_{2n}$ 满足条件：$EX_i=0$，$DX_i=1$，
$X_i$ 与 $X_j$（$i\ne j$）的相关系数为 $\rho$．
记 $\displaystyle\xi=\sum_{i=1}^{n}X_i$，$\displaystyle\eta=\sum_{j=n+1}^{2n}X_j$，
则 $\xi$ 与 $\eta$ 的相关系数 $\rho_{\xi\eta}$ 为
\begin{choices}[columns=4,label-pos=left]
\item $\dfrac{n\rho}{1+n\rho}$．
\item $n\rho$．
\item $\dfrac{n\rho}{1+(n+1)\rho}$．
\item $\dfrac{n\rho}{1+(n-1)\rho}$．
\end{choices}
\end{question}
\QuestionBreak

\begin{question}[index=9]
从总体 $X\sim N(\mu,2^2)$ 中随意抽取容量为 $n$ 的一组样本，其均值为 $\overline X$，
若要求 $E\bigl(\lvert\overline X-\mu\rvert\bigr)\le0.1$，则 $n$ 的最小取值为
\begin{choices}[columns=2]
\item $40$．
\item $128$．
\item $255$．
\item $1\,537$．
\end{choices}
\end{question}
\QuestionBreak

\begin{question}[index=10]
设两组数据分别来自两个方差相同且相互独立的正态总体 $X\sim N(\mu_1,\sigma^2)$
与 $Y\sim N(\mu_2,\sigma^2)$．现检验两组数据的均值有无显著差异，且犯第一类错误的概率为 $\alpha$，则应
\begin{choices}[columns=1]
\item 在置信度为 $1-\alpha$ 下，检验 $H_0:\ \mu_1\ne\mu_2$；$H_1:\ \mu_1=\mu_2$．
\item 在置信度为 $1-\dfrac{\alpha}{2}$ 下，检验 $H_0:\ \mu_1\ne\mu_2$；$H_1:\ \mu_1=\mu_2$．
\item 在置信度为 $1-\alpha$ 下，检验 $H_0:\ \mu_1=\mu_2$；$H_1:\ \mu_1\ne\mu_2$．
\item 在置信度为 $1-\dfrac{\alpha}{2}$ 下，检验 $H_0:\ \mu_1=\mu_2$；$H_1:\ \mu_1\ne\mu_2$．
\end{choices}
\end{question}
